\begin{helper}[A closed set separates outside two-color components]
Let $X$ be closed under a partial edge-coloring $\varphi$.  Suppose
$u\notin X$, $a,b\in X$, color $\alpha$ is missing at $a$, and color
$\beta$ is missing at $b$.  Then the full $(\alpha,\beta)$ incidence
component rooted at $u$ contains no vertex-node belonging to $X$ and no
edge-node incident with a vertex of $X$.
\end{helper}
\begin{proof}
Let $H$ be the two-color incidence graph from
\noderef{TwoColorIncidenceGraph}, and call a node of $H$ outside if it is
either a vertex-node $v$ with $v\notin X$, or an edge-node $e$ that is
incident with no vertex of $X$.  The root $u$ is outside by hypothesis.

We first check that every neighbor in $H$ of an outside node is outside.
There are only two nonvacuous cases because $H$ is bipartite.  If an
outside vertex-node $v$ is adjacent to an edge-node $e$, then $e$ is
incident with $v$ in the sense of \noderef{Incident}, and its color is
$\alpha$ or $\beta$.  If $e$ were also
incident with some $x\in X$, then $x\ne v$.  The two distinct incidences
identify the exact endpoint pair of $e$, so $e$ would belong to
\noderef{BoundaryEdges}.  In the $\alpha$ case, closedness as defined in
\noderef{ColoredVertexSetClosed}, applied at $a$, contradicts that
$\alpha$ is missing there; in the $\beta$ case the same argument at $b$
contradicts the missingness of $\beta$.  Thus $e$ is incident with no
vertex of $X$.  Conversely, if an outside edge-node $e$ is adjacent to a
vertex-node $v$, then $e$ is incident with $v$, so the defining property
of the edge-node immediately gives $v\notin X$.

Every node in the component of the root is joined to it by a finite walk,
by the connected-component definition in \noderef{TwoColorComponent}.
Induction along this walk, using the neighbor calculation above, shows
that every node in the component is outside.  Specializing first to
vertex-nodes and then to edge-nodes gives the two asserted conclusions.
\end{proof}
